\section{Fusão Sensorial Heterogênea}

A fusão sensorial heterogênea representa um campo fundamental na robótica autônoma, particularmente em ambientes complexos como subestações elétricas. Esta seção apresenta os conceitos essenciais, técnicas e desafios relacionados à integração de dados provenientes de múltiplos sensores de naturezas distintas para obtenção de uma percepção ambiental mais robusta e confiável.

\subsection{Conceitos Fundamentais de Fusão Sensorial}

A fusão sensorial pode ser definida como o processo de integração de dados provenientes de múltiplos sensores, com o objetivo de produzir informações mais precisas, completas e confiáveis do que seria possível utilizando cada sensor individualmente \cite{Hall1997}. Em robótica móvel, a percepção ambiental precisa ser robusta mesmo em condições adversas, como as encontradas em subestações elétricas \cite{Oliveira2021EstudoAplicacao}.

Em sistemas robóticos autônomos, a fusão sensorial busca superar as limitações inerentes a cada tipo de sensor, como limitações de alcance, sensibilidade a condições ambientais, ruído e incerteza nas medições, e interferência eletromagnética - particularmente relevante em subestações elétricas \cite{Liu2016EMC}.

A heterogeneidade na fusão sensorial refere-se à integração de dados provenientes de sensores de naturezas fundamentalmente distintas, como sensores visuais (câmeras \ac{RGB}), térmicos (câmeras termográficas), acústicos (microfones), de distância (\ac{LiDAR}, ultrassom), inerciais (\ac{IMU}) e químicos (sensores de gás) \cite{Kim2023MultimodalFusion}. Esta abordagem multimodal permite uma percepção mais completa do ambiente, explorando as complementaridades entre diferentes modalidades sensoriais.

\subsection{Níveis de Fusão Sensorial}

A fusão sensorial pode ser realizada em diferentes níveis de abstração, cada um com suas vantagens e desvantagens \cite{Costa2018ImplementacaoFusao, Santos2023FusaoSensorial}:

\subsubsection{Fusão de Baixo Nível (Nível de Dados)}

A fusão de baixo nível envolve a combinação direta dos sinais provenientes dos sensores antes de qualquer extração de características ou interpretação semântica:

\begin{equation}
Z_{fusão} = f(Z_1, Z_2, ..., Z_n)
\end{equation}

Onde $Z_i$ representa os dados brutos do sensor $i$, e $f$ é a função de fusão aplicada.

Exemplos incluem combinação de nuvens de pontos de múltiplos sensores \ac{LiDAR} e fusão de imagens de diferentes espectros (\ac{RGB} + infravermelho). A principal vantagem é a preservação de toda a informação original, mas requer maior capacidade computacional e pode ser sensível a problemas de sincronização temporal e calibração espacial entre sensores \cite{Schneider2018FusaoSensores}.

\subsubsection{Fusão de Médio Nível (Nível de Características)}

A fusão de médio nível opera sobre características extraídas dos dados sensoriais brutos:

\begin{equation}
F_{fusão} = g(F_1, F_2, ..., F_n)
\end{equation}

Onde $F_i$ representa o conjunto de características extraídas do sensor $i$, e $g$ é a função de fusão aplicada às características.

Este nível oferece um bom equilíbrio entre redução de dimensionalidade (e consequentemente de carga computacional) e preservação de informação relevante \cite{Grassi2006ArquiteturaHibrida}.

\subsubsection{Fusão de Alto Nível (Nível de Decisão)}

A fusão de alto nível combina decisões ou classificações já realizadas por algoritmos específicos para cada sensor:

\begin{equation}
D_{fusão} = h(D_1, D_2, ..., D_n)
\end{equation}

Onde $D_i$ representa a decisão ou classificação baseada no sensor $i$, e $h$ é a função de fusão aplicada às decisões.

Este nível requer menor largura de banda para comunicação entre subsistemas, mas pode perder informações importantes durante o processo de decisão individual de cada sensor \cite{Ramos2012AplicacaoTecnicas}.

\subsection{Técnicas de Fusão Sensorial}

\subsubsection{Filtro de Kalman e suas Variantes}

O Filtro de Kalman é uma das técnicas mais utilizadas para fusão sensorial, especialmente para estimação de estados em sistemas dinâmicos lineares com ruído gaussiano \cite{Oliveira2021EstudoAplicacao}. O filtro opera em dois estágios: predição e atualização.

Na etapa de predição, o estado atual é estimado com base no estado anterior e no modelo do sistema:

\begin{equation}
\hat{x}_{k|k-1} = F_k\hat{x}_{k-1|k-1} + B_ku_k
\end{equation}

\begin{equation}
P_{k|k-1} = F_kP_{k-1|k-1}F_k^T + Q_k
\end{equation}

Na etapa de atualização, a estimativa é corrigida com base nas medições dos sensores:

\begin{equation}
K_k = P_{k|k-1}H_k^T(H_kP_{k|k-1}H_k^T + R_k)^{-1}
\end{equation}

\begin{equation}
\hat{x}_{k|k} = \hat{x}_{k|k-1} + K_k(z_k - H_k\hat{x}_{k|k-1})
\end{equation}

Para sistemas não-lineares, variantes como o Filtro de Kalman Estendido (\ac{EKF}) e o Filtro de Kalman Unscented (\ac{UKF}) são utilizadas \cite{Costa2018ImplementacaoFusao}.

\subsubsection{Redes Bayesianas}

As redes Bayesianas oferecem um \textit{framework} probabilístico para modelar relações causais entre variáveis, sendo particularmente úteis para fusão sensorial de alto nível \cite{Ramos2012AplicacaoTecnicas}. A probabilidade conjunta de todas as variáveis na rede é dada por:

\begin{equation}
P(X_1, X_2, ..., X_n) = \prod_{i=1}^{n} P(X_i | Pa(X_i))
\end{equation}

Onde $Pa(X_i)$ representa o conjunto de pais do nó $X_i$ no grafo.

Em fusão sensorial, as redes Bayesianas permitem integrar informações de diferentes sensores considerando suas incertezas e relações causais.

\subsubsection{Redes Neurais e Aprendizado Profundo}

Técnicas de aprendizado de máquina, especialmente redes neurais profundas, têm sido cada vez mais aplicadas para fusão sensorial \cite{Kim2023MultimodalFusion}. Estas abordagens permitem aprender automaticamente representações e relações complexas entre diferentes modalidades sensoriais.

Arquiteturas comuns incluem Redes Neurais Convolucionais (\ac{CNN}) para processamento de dados espaciais, Redes Neurais Recorrentes (\ac{RNN}) para sequências temporais, e arquiteturas multimodais específicas para combinar diferentes tipos de dados sensoriais.

Uma arquitetura típica de fusão sensorial baseada em aprendizado profundo pode ser representada como:

\begin{equation}
F_{fusão} = f_{fusão}(f_1(X_1), f_2(X_2), ..., f_n(X_n))
\end{equation}

Onde $X_i$ representa os dados do sensor $i$, $f_i$ é a rede neural específica para processar o sensor $i$, e $f_{fusão}$ é a rede de fusão que combina as representações intermediárias.

\subsection{Desafios da Fusão em Ambientes com Interferência Eletromagnética}

A fusão sensorial em subestações elétricas enfrenta desafios específicos devido à presença de campos eletromagnéticos intensos \cite{Liu2016EMC, Wang2024CombinedNav}:

\begin{itemize}
    \item \textbf{Interferência em Sensores Eletrônicos}: Campos eletromagnéticos podem induzir ruídos e distorções em sensores eletrônicos, afetando a qualidade dos dados. Sensores magnéticos, inerciais e de comunicação sem fio são particularmente suscetíveis.
    
    \item \textbf{Variabilidade Temporal da Interferência}: A intensidade e distribuição espacial da interferência eletromagnética pode variar significativamente ao longo do tempo, exigindo sistemas de fusão sensorial adaptativos.
    
    \item \textbf{Heterogeneidade da Interferência}: Diferentes sensores são afetados de maneiras distintas pela interferência eletromagnética, o que deve ser considerado na estratégia de fusão.
\end{itemize}

Técnicas de mitigação incluem blindagem eletromagnética, filtragem digital e algoritmos de detecção e compensação de interferência \cite{Wang2024CombinedNav}.

\subsection{Modelos de Pesos e Bias para Relevância Sensorial}

Um aspecto crucial da fusão sensorial heterogênea é a atribuição de pesos relativos a cada sensor, refletindo sua confiabilidade, relevância e qualidade momentânea dos dados \cite{Oliveira2021EstudoAplicacao, Schneider2018FusaoSensores}.

\subsubsection{Ponderação Baseada em Covariância}

Um método comum para atribuir pesos aos sensores é baseado na matriz de covariância do erro de medição. Sensores com menor incerteza (menor covariância) recebem maior peso. Para o caso de fusão de estimativas escalares de múltiplos sensores, a estimativa ótima é dada por:

\begin{equation}
\hat{x} = \sum_{i=1}^{n} w_i x_i
\end{equation}

Onde os pesos $w_i$ são calculados como:

\begin{equation}
w_i = \frac{\sigma_i^{-2}}{\sum_{j=1}^{n} \sigma_j^{-2}}
\end{equation}

Com $\sigma_i^2$ representando a variância do erro de medição do sensor $i$.

\subsubsection{Ponderação Adaptativa Baseada em Contexto}

Em ambientes dinâmicos como subestações elétricas, a confiabilidade dos sensores pode variar em função do contexto operacional. Modelos adaptativos de ponderação consideram fatores contextuais para ajustar os pesos dos sensores em tempo real \cite{Wang2024CombinedNav}:

\begin{equation}
w_i(t) = f_w(w_i^{base}, c_1(t), c_2(t), ..., c_m(t))
\end{equation}

Onde $w_i(t)$ é o peso do sensor $i$ no instante $t$, $w_i^{base}$ é o peso base do sensor $i$, $c_j(t)$ são fatores contextuais (nível de interferência eletromagnética, iluminação, etc.), e $f_w$ é a função de adaptação dos pesos.

\subsection{Aplicações em Robótica Móvel para Inspeção}

A fusão sensorial heterogênea é particularmente relevante para robôs de inspeção em subestações elétricas, onde diferentes tipos de sensores são necessários para detectar diversos tipos de anomalias \cite{Lactec2022, Iberdrola2022}:

\begin{itemize}
    \item \textbf{Detecção de Pontos Quentes}: Fusão de imagens térmicas com visuais para identificação precisa de componentes superaquecidos.
    
    \item \textbf{Detecção de Descargas Parciais}: Combinação de sensores acústicos, \ac{UV} e \ac{RF} para localização de descargas parciais.
    
    \item \textbf{Navegação Robusta}: Fusão de \ac{LiDAR}, câmeras e \ac{IMU} para navegação precisa mesmo em condições de interferência eletromagnética.
    
    \item \textbf{Inspeção Visual Avançada}: Combinação de diferentes modalidades de imagem (visual, térmica, \ac{UV}) para detecção de anomalias visuais como corrosão, danos físicos ou vazamentos.
\end{itemize}

A integração eficaz destes sensores heterogêneos através de técnicas de fusão adequadas é fundamental para o desenvolvimento de sistemas robóticos autônomos capazes de realizar inspeções confiáveis e eficientes em subestações elétricas, superando os desafios impostos por este ambiente complexo.
